\documentclass{beamer}

\usepackage{amsmath, amsfonts, amssymb}
\usepackage{booktabs}
\usepackage{graphicx}
\usepackage{cmbright}
\usepackage{fontspec}

\beamertemplatenavigationsymbolsempty
\usetheme{metropolis}

\DeclareMathOperator*{\argmax}{arg\,max}

\AtBeginSection[]{
  \begin{frame}
  \vfill
  \centering
  \begin{beamercolorbox}[sep=8pt,center,shadow=true,rounded=true]{title}
    \usebeamerfont{title}\insertsectionhead\par%
  \end{beamercolorbox}
  \vfill
  \end{frame}
}

\title{Impulse Response Estimation by Smooth Local Projections (LP)}
\subtitle{Matthew Kielar \& Elijah Adrian}
\date[]{ECON 622 001}

\begin{document}

\begin{frame}
    \titlepage
\end{frame}

\subsection{Overview} 

\begin{frame}{Overview}

\end{frame}

\begin{frame}{Local Projections (LP)}

    Jorda (2005) proposed local projections (LP) as an alternative method for estimating impulse responses. \medskip

    Basic idea: Given variables $y_t$, instead of estimating one model just regress future values of
    $y$ on current values at each horizon $k$. \medskip
    
    At each horizon, run the following OLS regression:

    \begin{equation*}
        y_{t+k} = \alpha^{(k)} + \beta^{(k)}x_t + \Gamma^{(k)}\mathbf{W}_t + u_{t+k}
    \end{equation*}

    where $x_t$ are exogenous variable(s) and $W_t$ is a vector of control variables. 
   
    Under this setup, $\hat{\beta}^{(k)}$ captures the impulse response at horizon $k$. 

\end{frame}

\begin{frame}{Why LP?}

    VARs are efficient, but sensitive to misspecification. Incorrect specification at
    the first lag results in errors being iterated forward. \medskip

    Local projections are considerably more robust in this regard. \medskip 

    Estimation of separate OLS equations for each horizon $\to$ errors don't propagate
    across horizons. \medskip 

    \begin{itemize}

        \item Since IRF is just vector of scalar values, can apply different smoothing methods
        \item More flexible estimation options
        
    \end{itemize}

    Consistency: If the true DGP is VAR(p) and LP estimator includes the same lag structure, then
    $\hat{\Gamma}_1^{(k)} \to_p A_1^k$ as $T \to \infty$.

\end{frame}

\begin{frame}{Smooth Local Projections}
   
	Barnichon and Brownlees (2019) propose utilising B-splines to smooth the resulting sequence of LP estimates $\hat{\beta}^{(k)}$ \medskip

	IRFs estimated by LP can deviate from the true IRF $\to$ smoothing can bring the estimated impulse responses closer to the actual DGP \medskip

	B-splines are also a well-suited smoothing method here:

	\begin{itemize}
		
		\item Flexible and data-adaptive (easily fine-tuned via knots or polynomial degree)

		\item Computationally efficient due to local control properties 

		\item Can be regularised without too much hassle
	
	\end{itemize}

\end{frame}

\begin{frame}{B-Splines: An Overview}
    B-splines (basis splines) are locally supported polynomials smoothly connected at knots.
    \begin{itemize}
        \item Defined by the set of knots and degree of the polynomial
        \item Knots determine where the pieces of the B-spline meet, and the degree determines the smoothness of the resulting function
        \item Nonzero only on a small interval defined by the knots
    \end{itemize}
    A spline is constructed as a linear combination of these basis functions, with coefficients determined by the data. \medskip

    Hence we are estimating local projections as linear combinations of smoothly connected polynomials
    
\end{frame}

\begin{frame}{Implementation}

    Plan: write a generic python implementation of the smooth LP estimation procedure of Barnichon and Brownlees (2019) \medskip

    Objective: a fully-functioning set of methods to estimate IRFs by SLP given any arbitrary time-series dataframe \medskip	

    Some features we might consider (subject to revision): \medskip

    \begin{itemize}
        
        \item Method to handle exogenous variables (analogue to VAR-X)
        \item Additional smoothing options (kernel smoothers, etc.)
        \item Some plotting functionality
    
    \end{itemize}

\end{frame}

\end{document}